\subsection{Kernels, Images, Rank, and Nullity}
\begin{definition} \label{def:kernel-and-range-linear-map} \hlt{Kernel and Range}\\
For $T\in\LL(\mathcal{V},\mathcal{W})$, define
$\ker(T)\defeq\{\mathbf{v}\in\mathcal{V}:T\mathbf{v}=\mathbf{0}_{\mathcal{W}}\}$ and
$\operatorname{range}(T)\defeq\{T\mathbf{v}:\mathbf{v}\in\mathcal{V}\}$.\\
The kernel is also called the \hlt{nullspace} of $T$.
\end{definition}

\begin{definition} \hlt{Injective and Surjective Linear Maps}\\
A function $T:\mathcal{V}\to\mathcal{W}$ is \hlt{injective} if $T\mathbf{u}=T\mathbf{v}$ implies $\mathbf{u}=\mathbf{v}$. It is \hlt{surjective} if $\operatorname{range}(T)=\mathcal{W}$.
\end{definition}

\begin{proposition} \label{prop:kernel-and-range-are-subspaces} \hlt{Kernel and Range Are Subspaces}\\
If $T \in \LL(\mathcal{V}, \mathcal{W})$, then $\ker(T)$ is a subspace of $\mathcal{V}$ and $\operatorname{range}(T)$ is a subspace of $\mathcal{W}$.
\begin{proof}
Since $T\mathbf{0}_{\mathcal{V}}=\mathbf{0}_{\mathcal{W}}$, the kernel is nonempty.\\
If $\mathbf{u},\mathbf{v}\in\ker(T)$ and $a\in\mathbb{F}$, then $T(a\mathbf{u}+\mathbf{v})=aT\mathbf{u}+T\mathbf{v}=\mathbf{0}_{\mathcal{W}}$, so $a\mathbf{u}+\mathbf{v}\in\ker(T)$.\\
For the range, if $T\mathbf{u},T\mathbf{v}\in\operatorname{range}(T)$ and $a\in\mathbb{F}$, then $aT\mathbf{u}+T\mathbf{v}=T(a\mathbf{u}+\mathbf{v})\in\operatorname{range}(T)$.\\
\hlt{Structural Commentary}: The kernel and range are the two canonical subspaces attached to a linear map.\\
\hlt{Key Insight}: Linearity makes both preimages of zero and images of vectors closed under linear combinations.
\end{proof}
\end{proposition}

\begin{figure}[H]
\centering
\scalebox{0.75}{\begin{tikzpicture}[
  >=Stealth,
  every node/.style={font=\small},
  ambient/.style={draw=gray!35, line width=0.45pt},
  subspace/.style={draw=gray!60, fill=gray!6, line width=0.45pt},
  maparrow/.style={->, gray!65, line width=0.45pt},
  point/.style={circle, fill=gray!70, inner sep=1.4pt}
]
  \coordinate (vzero) at (2.0,1.75);
  \coordinate (wzero) at (7.55,1.75);

  \draw[ambient, fill=white] (0,0) rectangle (4.2,3.6);
  \draw[ambient, fill=gray!15] (5.6,0) rectangle (9.8,3.6);

  \node[anchor=north west] at (0.28,3.42) {$\mathcal{V}$};
  \node[anchor=north west] at (5.88,3.42) {$\mathcal{W}$};

  \draw[subspace]
    plot[smooth cycle, tension=0.75] coordinates {
      (0.78,1.20)
      (0.83,1.85)
      (1.24,2.58)
      (1.98,2.70)
      (2.80,2.25)
      (3.38,1.80)
      (3.02,1.10)
      (2.38,0.76)
      (1.43,0.86)
    };
  \node at (1.85,2.32) {$\ker(T)$};

  \draw[subspace, fill=gray!4]
    plot[smooth cycle, tension=0.72] coordinates {
      (6.24,1.05)
      (6.45,1.72)
      (6.95,2.55)
      (7.60,3.10)
      (8.25,3.02)
      (8.78,2.33)
      (8.94,1.54)
      (8.52,0.90)
      (7.62,0.38)
      (6.72,0.55)
    };
  \node at (7.63,2.34) {$\operatorname{range}(T)$};

  \node[point] at (vzero) {};
  \node[below=4pt] at (vzero) {$\mathbf{0}$};
  \node[point] at (wzero) {};
  \node[below=4pt] at (wzero) {$\mathbf{0}$};

  \draw[maparrow] (vzero) -- (wzero);
  \draw[maparrow] (4.2,3.6) -- (6.62,2.58);
  \draw[maparrow] (4.2,0.0) -- (6.55,1.03);
\end{tikzpicture}
}
\caption{For a linear map $T:\mathcal{V}\to\mathcal{W}$, the range of $T$ may be a proper subspace of $\mathcal{W}$, and the kernel may contain nonzero vectors. The conditions $\operatorname{range}(T)=\mathcal{W}$ and $\ker(T)=\{\mathbf{0}\}$ characterize surjectivity and injectivity, respectively.}
\label{fig:kernel-range-linear-map}
\end{figure}

\begin{definition} \label{def:projection-onto-subspace-along-complement} \hlt{Projection Along a Complement}\\
An operator $P\in\LL(\mathcal{V})$ is an \hlt{algebraic projection}, or \hlt{idempotent operator}, if $P^2=P$.\\
If $\mathcal{V}=\mathcal{U}\oplus\mathcal{W}$, the \hlt{projection onto $\mathcal{U}$ along $\mathcal{W}$} is $\rho_{\mathcal{U},\mathcal{W}}\in\LL(\mathcal{V})$ where $\rho_{\mathcal{U},\mathcal{W}}(\mathbf{u}+\mathbf{w})\defeq\mathbf{u}$ for $\mathbf{u}\in\mathcal{U}$, $\mathbf{w}\in\mathcal{W}$.
\end{definition}

\begin{proposition} \label{prop:algebraic-projection-range-kernel-idempotent} \hlt{Algebraic Projections and Direct-Sum Splittings}\\
For $P\in\LL(\mathcal{V})$, $P$ is an algebraic projection if and only if $\mathcal{V}=\operatorname{range}(P)\oplus\ker(P)$ and $P=\rho_{\operatorname{range}(P),\ker(P)}$.\\
In particular, if $\mathcal{V}=\mathcal{U}\oplus\mathcal{W}$ and $\rho=\rho_{\mathcal{U},\mathcal{W}}$, then $\rho$ is an algebraic projection, $\operatorname{range}(\rho)=\mathcal{U}$, and $\ker(\rho)=\mathcal{W}$.
\begin{proof}
\begin{enumerate}[label=\roman*.]
\setlength{\itemsep}{0pt}
\item[$(\Rightarrow)$] Suppose $P$ is an algebraic projection, so $P^2=P$. For every $\mathbf{v}\in\mathcal{V}$, $\mathbf{v}=P\mathbf{v}+(\mathbf{v}-P\mathbf{v})$, with $P\mathbf{v}\in\operatorname{range}(P)$ and $\mathbf{v}-P\mathbf{v}\in\ker(P)$. If $\mathbf{x}\in\operatorname{range}(P)\cap\ker(P)$, then $\mathbf{x}=P\mathbf{y}$ and $\mathbf{0}=P\mathbf{x}=P^2\mathbf{y}=\mathbf{x}$. Thus the sum is direct. Also $P$ fixes $\operatorname{range}(P)$ and kills $\ker(P)$, so $P=\rho_{\operatorname{range}(P),\ker(P)}$.
\item[$(\Leftarrow)$] Assume the right-hand side. For $\mathbf{v}=\mathbf{x}+\mathbf{y}$ with $\mathbf{x}\in\operatorname{range}(P)$ and $\mathbf{y}\in\ker(P)$, $P^2\mathbf{v}=P\mathbf{x}=\mathbf{x}=P\mathbf{v}$.\\
Thus $P$ is an algebraic projection. For $\rho_{\mathcal{U},\mathcal{W}}$, the formula $\rho_{\mathcal{U},\mathcal{W}}(\mathbf{u}+\mathbf{w})=\mathbf{u}$ gives range $\mathcal{U}$ and kernel $\mathcal{W}$, and the equivalence above shows that $\rho$ is an algebraic projection.
\end{enumerate}
\hlt{Structural Commentary}: Idempotents keep one summand and kill its complement.\\
\hlt{Key Insight}: Direct-sum coordinates turn a projection into $I\oplus0$.
\end{proof}
\end{proposition}

\begin{proposition} \label{prop:complementary-algebraic-projections} \hlt{Complementary Algebraic Projections}\\
If $\mathcal{V}=\mathcal{U}\oplus\mathcal{W}$, then $\rho_{\mathcal{U},\mathcal{W}}+\rho_{\mathcal{W},\mathcal{U}}=I_{\mathcal{V}}$, and $\rho_{\mathcal{U},\mathcal{W}}\rho_{\mathcal{W},\mathcal{U}} =\rho_{\mathcal{W},\mathcal{U}}\rho_{\mathcal{U},\mathcal{W}}=0$.
\begin{proof}
For $\mathbf{v}=\mathbf{u}+\mathbf{w}$, $\mathbf{u}\in\mathcal{U}$, $\mathbf{w}\in\mathcal{W}$, $(\rho_{\mathcal{U},\mathcal{W}}+\rho_{\mathcal{W},\mathcal{U}})\mathbf{v} =\mathbf{u}+\mathbf{w} =\mathbf{v}$,\\
and $\rho_{\mathcal{U},\mathcal{W}}\rho_{\mathcal{W},\mathcal{U}}\mathbf{v} =\rho_{\mathcal{U},\mathcal{W}}\mathbf{w} =\mathbf{0}$; $\rho_{\mathcal{W},\mathcal{U}}\rho_{\mathcal{U},\mathcal{W}}\mathbf{v} =\rho_{\mathcal{W},\mathcal{U}}\mathbf{u} =\mathbf{0}$. Thus the stated operator identities hold.\\
\hlt{Structural Commentary}: A direct-sum decomposition gives two complementary idempotents.\\
\hlt{Key Insight}: Projection onto one summand kills the other summand.
\end{proof}
\end{proposition}

\begin{definition} \label{def:algebraically-orthogonal-projections-and-resolutions} \hlt{Orthogonal Projections in the Algebraic Sense}\\
Two algebraic projections $\rho,\sigma\in\LL(\mathcal{V})$ are \hlt{algebraically orthogonal}, written $\rho\perp\sigma$, if $\rho\sigma=\sigma\rho=0$.\\
\hlt{Resolution of identity} on $\mathcal{V}$ is a finite family of algebraic projections $(\rho_i)_{i=1}^{k}$ with $\rho_i\perp\rho_j$ for $i\neq j$ and $\sum\limits_{i=1}^{k}\rho_i=I_{\mathcal{V}}$.
\end{definition}

\begin{theorem} \label{thm:resolution-of-identity-gives-direct-sum} \hlt{Resolution Ranges Form a Direct Sum}\\
If $(\rho_i)_{i=1}^{k}$ is a resolution of $I_{\mathcal{V}}$ and $R_i=\operatorname{range}(\rho_i)$, then $\mathcal{V}=\bigoplus\limits_{i=1}^{k}R_i$, and $\ker(\rho_i)=\bigoplus\limits_{j\neq i}R_j$.
\begin{proof}
Since $\sum\limits_{i=1}^{k}\rho_i=I_{\mathcal{V}}$, $\mathbf{v}=\sum\limits_{i=1}^{k}\rho_i\mathbf{v}\in\sum\limits_{i=1}^{k}R_i$.\\
If $\sum\limits_{i=1}^{k}\mathbf{x}_i=\mathbf{0}$ with $\mathbf{x}_i\in R_i$, then for each $1\leq m\leq k$, $\mathbf{0}=\rho_m\sum\limits_{i=1}^{k}\mathbf{x}_i=\mathbf{x}_m$, as $\rho_m$ fixes $R_m$ and kills $R_i$ for $i\neq m$.\\
Thus $\mathcal{V}=\bigoplus\limits_{i=1}^{k}R_i$. Also, for $\mathbf{v}=\sum\limits_{j=1}^{k}\mathbf{x}_j$ with $\mathbf{x}_j\in R_j$, $\rho_i\mathbf{v}=\mathbf{x}_i$, so $\ker(\rho_i)=\bigoplus\limits_{j\neq i}R_j$.\\
\hlt{Structural Commentary}: A resolution decomposes the identity into independent idempotent pieces.\\
\hlt{Key Insight}: Pairwise zero products are the operator form of directness.
\end{proof}
\end{theorem}

\begin{theorem} \label{thm:direct-sum-gives-resolution-of-identity} \hlt{Direct Sums Give Resolutions}\\
If $\mathcal{V}=\bigoplus_{i=1}^{k}\mathcal{U}_i$ and $\rho_i=\rho_{\mathcal{U}_i,\bigoplus_{j\neq i}\mathcal{U}_j}$, then $(\rho_i)_{i=1}^{k}$ is a resolution of $I_{\mathcal{V}}$.
\begin{proof}
For $\mathbf{v}=\sum_{i=1}^{k}\mathbf{u}_i$, $\mathbf{u}_i\in\mathcal{U}_i$, $\sum\limits_{i=1}^{k}\rho_i\mathbf{v} =\sum\limits_{i=1}^{k}\mathbf{u}_i =\mathbf{v}$. Thus $\sum\limits_{i=1}^{k}\rho_i=I_{\mathcal{V}}$.\\
If $i\neq j$, then $\rho_i\rho_j\mathbf{v}=\rho_i\mathbf{u}_j=\mathbf{0}$ and $\rho_j\rho_i\mathbf{v}=\rho_j\mathbf{u}_i=\mathbf{0}$, so $\rho_i\perp\rho_j$.\\
\hlt{Structural Commentary}: Direct-sum coordinates produce mutually annihilating coordinate projections.\\
\hlt{Key Insight}: Coordinate projections add to the identity and kill all other summands.
\end{proof}
\end{theorem}

\begin{theorem} \label{thm:algebra-of-projections} \hlt{Algebra of Projections}\\
Let $\rho,\sigma\in\LL(\mathcal{V})$ be algebraic projections and $\operatorname{char}(\mathbb{F})\neq2$.
\begin{enumerate}[label=(\roman*)]
\setlength{\itemsep}{0pt}
\item $\rho+\sigma$ is an algebraic projection $\Longleftrightarrow \rho\perp\sigma$.\\
Then $\operatorname{range}(\rho+\sigma)=\operatorname{range}(\rho)\oplus\operatorname{range}(\sigma)$; $\ker(\rho+\sigma)=\ker(\rho)\cap\ker(\sigma)$
\item $\rho-\sigma$ is an algebraic projection $\Longleftrightarrow \rho\sigma=\sigma\rho=\sigma$.\\
Then $\operatorname{range}(\rho-\sigma)=\operatorname{range}(\rho)\cap\ker(\sigma)$; $\ker(\rho-\sigma)=\ker(\rho)\oplus\operatorname{range}(\sigma)$.
\item If $\rho\sigma=\sigma\rho$, then $\rho\sigma$ is algebraic projection and $\operatorname{range}(\rho\sigma)=\operatorname{range}(\rho)\cap\operatorname{range}(\sigma)$; $\ker(\rho\sigma)=\ker(\rho)+\ker(\sigma)$.
\end{enumerate}
\begin{proof}
\begin{enumerate}[label=(\roman*)]
\setlength{\itemsep}{0pt}
\item Put $A=\rho\sigma$, $B=\sigma\rho$, and $C=\rho\sigma\rho$. Since $(\rho+\sigma)^2=\rho+\sigma \Longleftrightarrow A+B=0$,\\
$A+B=0$ gives $A+C=\rho(A+B)=0$ and $C+B=(A+B)\rho=0$, hence $B=A$ and $2A=0$.\\
Thus $A=B=0$. The converse is immediate. Assume $\rho\perp\sigma$. Then $\operatorname{range}(\rho+\sigma)\subseteq\operatorname{range}(\rho)+\operatorname{range}(\sigma)$, while $(\rho+\sigma)(\mathbf{x}+\mathbf{y})=\mathbf{x}+\mathbf{y}$ for $\mathbf{x}\in\operatorname{range}(\rho)$, $\mathbf{y}\in\operatorname{range}(\sigma)$. Also $\operatorname{range}(\rho)\cap\operatorname{range}(\sigma)=\{\mathbf{0}\}$, since $\mathbf{z}\in\operatorname{range}(\rho)\cap\operatorname{range}(\sigma)$ and $\mathbf{z}=\sigma\mathbf{w}$ imply $\mathbf{z}=\rho\mathbf{z}=\rho\sigma\mathbf{w}=\mathbf{0}$. Finally, $(\rho+\sigma)\mathbf{v}=0 \Longleftrightarrow \rho\mathbf{v}=0=\sigma\mathbf{v}$.
\item Since $P^2=P\Longleftrightarrow (I_{\mathcal{V}}-P)^2=I_{\mathcal{V}}-P$, $\rho-\sigma$ is an algebraic projection $\Longleftrightarrow$ $(I_{\mathcal{V}}-\rho)+\sigma$ is an algebraic projection. By (i), this is equivalent to $(I_{\mathcal{V}}-\rho)\sigma=\sigma(I_{\mathcal{V}}-\rho)=0$, i.e. $\rho\sigma=\sigma\rho=\sigma$.\\
Assume $\rho\sigma=\sigma\rho=\sigma$ and set $Q=\rho-\sigma$. Then $Q\mathbf{v}\in\operatorname{range}(\rho)\cap\ker(\sigma)$; $Q\mathbf{x}=\mathbf{x}$, $\mathbf{x}\in\operatorname{range}(\rho)\cap\ker(\sigma)$, so $\operatorname{range}(Q)=\operatorname{range}(\rho)\cap\ker(\sigma)$. Also $Q$ kills both $\ker(\rho)$ and $\operatorname{range}(\sigma)$, and $\ker(\rho)\cap\operatorname{range}(\sigma)=\{\mathbf{0}\}$.\\
If $\mathbf{v}\in\ker(Q)$, then $\rho\mathbf{v}=\sigma\mathbf{v}$ and $\mathbf{v}=(\mathbf{v}-\rho\mathbf{v})+\sigma\mathbf{v} \in\ker(\rho)\oplus\operatorname{range}(\sigma)$. Hence $\ker(Q)=\ker(\rho)\oplus\operatorname{range}(\sigma)$.
\item If $\rho\sigma=\sigma\rho$, then $(\rho\sigma)^2=\rho^2\sigma^2=\rho\sigma$. Moreover $\mathbf{x}\in\operatorname{range}(\rho\sigma)$ $\Longleftrightarrow$ $\rho\mathbf{x}=\mathbf{x}=\sigma\mathbf{x}$,\\
so $\operatorname{range}(\rho\sigma)=\operatorname{range}(\rho)\cap\operatorname{range}(\sigma)$. Finally, $\ker(\rho)+\ker(\sigma)\subseteq\ker(\rho\sigma)$, and if $\rho\sigma\mathbf{v}=0$,\\
then $\mathbf{v}=(\mathbf{v}-\sigma\mathbf{v})+\sigma\mathbf{v}$; $\mathbf{v}-\sigma\mathbf{v}\in\ker(\sigma)$, $\sigma\mathbf{v}\in\ker(\rho)$. Thus $\ker(\rho\sigma)=\ker(\rho)+\ker(\sigma)$.
\end{enumerate}
\hlt{Structural Commentary}: Operator identities encode range-kernel geometry.\\
\hlt{Key Insight}: Idempotence is a rigid equation; expanding it forces the cross terms to encode the geometry.
\end{proof}
\end{theorem}

\begin{proposition} \label{prop:shift-operators-finitely-nonzero} \hlt{Shift Operators on Finitely Nonzero Sequences}\\
Define $R,S:\mathbb{F}^{(\mathbb{N})}\to\mathbb{F}^{(\mathbb{N})}$ by $R(x_1,x_2,\ldots)=(0,x_1,x_2,\ldots)$, $S(x_1,x_2,\ldots)=(x_2,x_3,\ldots)$. Then $R,S\in\LL(\mathbb{F}^{(\mathbb{N})})$.\\
Also, $\ker(R)=\{\mathbf{0}\}$, $\operatorname{range}(R)=\{(x_1,x_2,\ldots)\in\mathbb{F}^{(\mathbb{N})}:x_1=0\}$, $\ker(S)=\operatorname{span}((1,0,0,\ldots))$, $\operatorname{range}(S)=\mathbb{F}^{(\mathbb{N})}$.
\begin{proof}
Both $R$ and $S$ preserve finite support, so they map $\mathbb{F}^{(\mathbb{N})}$ into itself.\\
Since both maps act componentwise by shifting entries, $R(a\mathbf{x}+\mathbf{y})=aR\mathbf{x}+R\mathbf{y}$ and $S(a\mathbf{x}+\mathbf{y})=aS\mathbf{x}+S\mathbf{y}$ for all $\mathbf{x},\mathbf{y}\in\mathbb{F}^{(\mathbb{N})}$ and $a\in\mathbb{F}$. Thus $R$ and $S$ are linear.\\
The map $R$ has zero kernel as $R\mathbf{x}=\mathbf{0}$ forces every entry of $\mathbf{x}$ to be zero, and its range is exactly the set of finitely nonzero sequences with first entry zero. For $S$, $S\mathbf{x}=\mathbf{0}$ forces $x_2=x_3=\cdots=0$, so $\ker(S)=\operatorname{span}((1,0,0,\ldots))$. Finally, for any $(y_1,y_2,\ldots)\in\mathbb{F}^{(\mathbb{N})}$, we have $S(0,y_1,y_2,\ldots)=(y_1,y_2,\ldots)$, so $S$ is surjective.\\
\hlt{Structural Commentary}: Shift operators expose behavior that finite-dimensional maps cannot have.\\
\hlt{Key Insight}: The right shift is injective but not surjective, while the left shift is surjective but not injective.
\end{proof}
\end{proposition}

\begin{example} \label{ex:shift-operators-show-operator-algebra-noncommutative} \hlt{Shift Operators Show Noncommutativity}\\
For the shifts $R,S\in\LL(\mathbb{F}^{(\mathbb{N})})$ from Proposition \ref{prop:shift-operators-finitely-nonzero}, $SR=I_{\mathbb{F}^{(\mathbb{N})}}$ and $RS(x_1,x_2,\ldots)=(0,x_2,x_3,\ldots)$.\\
Hence $RS\neq SR$, for instance on $(1,0,0,\ldots)$. Therefore $\LL(\mathbb{F}^{(\mathbb{N})})$ is not commutative under composition.
\end{example}

\begin{proposition} \label{prop:injective-iff-kernel-zero} \hlt{Injectivity and the Kernel}\\
A linear map $T \in \LL(\mathcal{V}, \mathcal{W})$ is injective if and only if $\ker(T) = \{ \mathbf{0} \}$.
\begin{proof}
\begin{enumerate}[label=\roman*.]
\setlength{\itemsep}{0pt}
\item[$(\Rightarrow)$] If $T$ is injective and $\mathbf{v}\in\ker(T)$, then $T\mathbf{v}=\mathbf{0}=T\mathbf{0}$, so $\mathbf{v}=\mathbf{0}$.
\item[$(\Leftarrow)$] If $\ker(T)=\{\mathbf{0}\}$ and $T\mathbf{v}=T\mathbf{w}$, then $T(\mathbf{v}-\mathbf{w})=\mathbf{0}$. Hence $\mathbf{v}-\mathbf{w}\in\ker(T)$, so $\mathbf{v}=\mathbf{w}$.
\end{enumerate}
\hlt{Structural Commentary}: The kernel measures exactly where a linear map collapses vectors to zero.\\
\hlt{Key Insight}: Injectivity is equivalent to collapsing only the zero vector.
\end{proof}
\end{proposition}

\begin{theorem} \label{thm:range-spanned-by-images-of-basis} \hlt{Range from Basis Images}\\
Let $\mathcal{V}$ be finite-dimensional with basis $(\mathbf{v}_1,\ldots,\mathbf{v}_n)$, and let $T\in\LL(\mathcal{V},\mathcal{W})$. Then $\operatorname{range}(T)=\operatorname{span}(T\mathbf{v}_1,\ldots,T\mathbf{v}_n)$.\\
In particular, $\operatorname{range}(T)$ is finite-dimensional and $\dim(\operatorname{range}(T))\leq n$.
\begin{proof}
If $\mathbf{v}=c_1\mathbf{v}_1+\cdots+c_n\mathbf{v}_n$, then $T\mathbf{v}=c_1T\mathbf{v}_1+\cdots+c_nT\mathbf{v}_n$.\\
Thus every vector in $\operatorname{range}(T)$ lies in $\operatorname{span}(T\mathbf{v}_1,\ldots,T\mathbf{v}_n)$, and the reverse inclusion is immediate from the definition of range. The dimension bound follows because this spanning list has length $n$.\\
\hlt{Structural Commentary}: The range is generated by the images of any basis of the domain.\\
\hlt{Key Insight}: Linear maps carry a basis-spanning formula directly into a range-spanning formula.
\end{proof}
\end{theorem}

\begin{definition} \label{def:rank-and-nullity-linear-map} \hlt{Rank and Nullity}\\
For a finite-dimensional domain, the rank of $T$ is $\dim(\operatorname{range}(T))$ and the nullity of $T$ is $\dim(\ker(T))$.
\end{definition}

\begin{theorem} \label{thm:rank-nullity} \hlt{Rank-Nullity Theorem}\\
If $\mathcal{V}$ is finite-dim and $T \in \LL(\mathcal{V}, \mathcal{W})$, then $\operatorname{range}(T)$ is finite-dim and $\dim(\mathcal{V}) = \dim(\ker(T)) + \dim(\operatorname{range}(T))$.
\begin{proof}
By Proposition \ref{prop:kernel-and-range-are-subspaces} and Corollary \ref{cor:finite-dimensional-spaces-and-subspaces-have-bases}, $\ker(T)$ has a basis $(\mathbf{u}_1,\ldots,\mathbf{u}_m)$.\\
Extend it to a basis $(\mathbf{u}_1,\ldots,\mathbf{u}_m,\mathbf{w}_1,\ldots,\mathbf{w}_r)$ of $\mathcal{V}$ by Theorem \ref{thm:extending-independent-lists}.\\
If $\mathbf{v}=\sum\limits_{i=1}^m a_i\mathbf{u}_i+\sum\limits_{j=1}^r b_j\mathbf{w}_j$, then $T\mathbf{v}=\sum\limits_{j=1}^r b_jT\mathbf{w}_j$, because $T\mathbf{u}_i=\mathbf{0}$. Hence $(T\mathbf{w}_1,\ldots,T\mathbf{w}_r)$ spans $\operatorname{range}(T)$.\\
If $\sum\limits_{j=1}^r c_jT\mathbf{w}_j=\mathbf{0}$, then $\sum\limits_{j=1}^r c_j\mathbf{w}_j\in\ker(T)$, so for some scalars $d_i$, $\sum\limits_{j=1}^r c_j\mathbf{w}_j=\sum\limits_{i=1}^m d_i\mathbf{u}_i$.\\
Thus $\sum\limits_{j=1}^r c_j\mathbf{w}_j-\sum\limits_{i=1}^m d_i\mathbf{u}_i=\mathbf{0}$, and linear independence of extended basis forces every $c_j$ and $d_i$ to be zero.\\
Hence $(T\mathbf{w}_1,\ldots,T\mathbf{w}_r)$ is a basis of $\operatorname{range}(T)$. Thus $\dim(\ker(T))=m$, $\dim(\operatorname{range}(T))=r$, and $\dim(\mathcal{V})=m+r$.\\
\hlt{Structural Commentary}: An adapted basis splits $\mathcal{V}$ into kernel directions and directions detected by $T$.\\
\hlt{Key Insight}: Rank-nullity says every domain dimension is either killed by $T$ or appears in its range.
\end{proof}
\end{theorem}

\begin{corollary} \label{cor:map-to-smaller-space-not-injective} \hlt{Maps to Smaller Spaces Are Not Injective}\\
Let $\mathcal{V}$ and $\mathcal{W}$ be finite-dimensional vector spaces with $\dim(\mathcal{V})>\dim(\mathcal{W})$.\\
Then no linear map from $\mathcal{V}$ to $\mathcal{W}$ is injective.
\begin{proof}
Let $T\in\LL(\mathcal{V},\mathcal{W})$. By Theorem \ref{thm:rank-nullity}, $\dim(\ker(T))=\dim(\mathcal{V})-\dim(\operatorname{range}(T)) \geq \dim(\mathcal{V})-\dim(\mathcal{W})>0$.\\
Thus $\ker(T)$ contains a nonzero vector, so $T$ is not injective by Proposition \ref{prop:injective-iff-kernel-zero}.\\
\hlt{Structural Commentary}: Excess domain dimension forces collapse.\\
\hlt{Key Insight}: The dimension deficit must appear inside the kernel.
\end{proof}
\end{corollary}

\begin{corollary} \label{cor:map-to-larger-space-not-surjective} \hlt{Maps to Larger Spaces Are Not Surjective}\\
Let $\mathcal{V}$ and $\mathcal{W}$ be finite-dimensional vector spaces with $\dim(\mathcal{V})<\dim(\mathcal{W})$.\\
Then no linear map from $\mathcal{V}$ to $\mathcal{W}$ is surjective.
\begin{proof}
Let $T\in\LL(\mathcal{V},\mathcal{W})$. By Theorem \ref{thm:rank-nullity}, $\dim(\operatorname{range}(T))=\dim(\mathcal{V})-\dim(\ker(T))\leq\dim(\mathcal{V})<\dim(\mathcal{W})$.\\
Hence $\operatorname{range}(T)$ is a proper subspace of $\mathcal{W}$, so $T$ is not surjective.\\
\hlt{Structural Commentary}: A smaller domain cannot cover a larger codomain.\\
\hlt{Key Insight}: The range can never have dimension larger than the domain.
\end{proof}
\end{corollary}

\begin{corollary} \label{cor:equal-dimension-kernel-range} \hlt{Equal-Dimension Kernel-Range Criterion}\\
Let $\mathcal{V}$ and $\mathcal{W}$ be finite-dimensional vector spaces with $\dim(\mathcal{V})=\dim(\mathcal{W})$, and let $T\in\LL(\mathcal{V},\mathcal{W})$.\\
Then $\ker(T)=\{\mathbf{0}\}$ if and only if $\operatorname{range}(T)=\mathcal{W}$. Equivalently, $T$ is injective if and only if $T$ is surjective.
\begin{proof}
\begin{enumerate}[label=\roman*.]
\setlength{\itemsep}{0pt}
\item[$(\Rightarrow)$] If $\ker(T)=\{\mathbf{0}\}$, then rank-nullity gives $\dim(\operatorname{range}(T))=\dim(\mathcal{V})=\dim(\mathcal{W})$.\\
Since $\operatorname{range}(T)$ is a subspace of $\mathcal{W}$ with the same dimension, $\operatorname{range}(T)=\mathcal{W}$.
\item[$(\Leftarrow)$] If $\operatorname{range}(T)=\mathcal{W}$, then rank-nullity gives $\dim(\ker(T))=\dim(\mathcal{V})-\dim(\mathcal{W})=0$, hence $\ker(T)=\{\mathbf{0}\}$.
\end{enumerate}
\hlt{Structural Commentary}: Equal finite dimensions make injectivity and surjectivity equivalent.\\
\hlt{Key Insight}: In equal dimensions, having no kernel loss is exactly having full range.
\end{proof}
\end{corollary}

\begin{definition} \hlt{Invertible Linear Map and Isomorphism}\\
A linear map $T \in \LL(\mathcal{V}, \mathcal{W})$ is \hlt{invertible} if there exists $S \in \LL(\mathcal{W}, \mathcal{V})$ such that $ST = I_{\mathcal{V}}$ and $TS = I_{\mathcal{W}}$.\\
An invertible linear map is called an \hlt{isomorphism}, and if one exists we write $\mathcal{V} \cong \mathcal{W}$.
\end{definition}

\begin{proposition} \label{prop:inverse-linear-map-is-unique} \hlt{Inverse Is Unique}\\
An invertible linear map has a unique inverse. Thus if $T$ is invertible, its inverse is denoted $T^{-1}$.
\begin{proof}
Suppose $S_1,S_2\in\LL(\mathcal{W},\mathcal{V})$ are inverses of $T\in\LL(\mathcal{V},\mathcal{W})$.\\
Then $S_1 = S_1I_{\mathcal{W}} = S_1(TS_2) = (S_1T)S_2 = I_{\mathcal{V}}S_2 = S_2$. Hence the inverse is unique.\\
\hlt{Structural Commentary}: A two-sided inverse is forced by the composition identities.\\
\hlt{Key Insight}: Associativity lets the two candidate inverses meet around $T$ and collapse to the same map.
\end{proof}
\end{proposition}

\begin{proposition} \label{prop:inverse-of-product-linear-maps} \hlt{Inverse of a Product Reverses Order}\\
Let $T\in\LL(\mathcal{U},\mathcal{V})$ and $S\in\LL(\mathcal{V},\mathcal{W})$ be invertible. Then $ST$ is invertible and $(ST)^{-1}=T^{-1}S^{-1}$.
\begin{proof}
Since $T^{-1}\in\LL(\mathcal{V},\mathcal{U})$ and $S^{-1}\in\LL(\mathcal{W},\mathcal{V})$, the product $T^{-1}S^{-1}\in\LL(\mathcal{W},\mathcal{U})$ is defined.\\
Moreover, $(T^{-1}S^{-1})(ST)=T^{-1}(S^{-1}S)T=T^{-1}T=I_{\mathcal{U}}$ and $(ST)(T^{-1}S^{-1})=S(TT^{-1})S^{-1}=SS^{-1}=I_{\mathcal{W}}$.\\
Thus $T^{-1}S^{-1}$ is a two-sided inverse of $ST$, so $(ST)^{-1}=T^{-1}S^{-1}$.\\
\hlt{Structural Commentary}: To undo a composition, undo the last map first.\\
\hlt{Key Insight}: Composition order reverses under inversion.
\end{proof}
\end{proposition}

\begin{proposition} \label{prop:linear-map-invertible-iff-bijective} \hlt{Invertibility and Bijectivity}\\
A linear map is invertible if and only if it is both injective and surjective.
\begin{proof}
\begin{enumerate}[label=\roman*.]
\setlength{\itemsep}{0pt}
\item[$(\Rightarrow)$] Let $S$ be an inverse of $T$. If $T\mathbf{v}=T\mathbf{w}$, then $\mathbf{v}=ST\mathbf{v}=ST\mathbf{w}=\mathbf{w}$, so $T$ is injective. Also every $\mathbf{w}$ in the codomain satisfies $\mathbf{w}=T(S\mathbf{w})$, so $T$ is surjective.
\item[$(\Leftarrow)$] Suppose $T$ is injective and surjective. Define $S\mathbf{w}$ to be the unique vector $\mathbf{v}$ with $T\mathbf{v}=\mathbf{w}$.\\
Since both sides below map under $T$ to $a\mathbf{w}_1+\mathbf{w}_2$, $S(a\mathbf{w}_1+\mathbf{w}_2)=aS\mathbf{w}_1+S\mathbf{w}_2$. Thus $S$ is linear and satisfies $ST=I$ and $TS=I$, so $T$ is invertible.
\end{enumerate}
\hlt{Structural Commentary}: Bijectivity gives a unique reverse map, and linearity makes it a linear inverse.\\
\hlt{Key Insight}: A linear map is invertible exactly when it is one-to-one and onto.
\end{proof}
\end{proposition}

\begin{definition} \label{def:linear-map-morphism-terminology-vector-spaces} \hlt{Morphism Terminology for Linear Maps}\\
For $T\in\LL(\mathcal{V},\mathcal{W})$:
\begin{enumerate}[label=(\roman*)]
\setlength{\itemsep}{0pt}
\item $T$ is a \hlt{homomorphism}.
\item If $\mathcal{V}=\mathcal{W}$, then $T$ is an \hlt{endomorphism}.
\item If $\ker(T)=\{\mathbf{0}\}$, then $T$ is a \hlt{monomorphism}, or \hlt{embedding}.
\item If $\operatorname{range}(T)=\mathcal{W}$, then $T$ is an \hlt{epimorphism}.
\item If $\ker(T)=\{\mathbf{0}\}$ and $\operatorname{range}(T)=\mathcal{W}$, then $T$ is an \hlt{isomorphism}.
\item If $\mathcal{V}=\mathcal{W}$ and $T$ is invertible, then $T$ is an \hlt{automorphism}.
\end{enumerate}
\end{definition}

\begin{proposition} \label{prop:rank-nullity-complement-view} \hlt{Rank-Nullity as a Complement Isomorphism}\\
Let $T\in\LL(\mathcal{V},\mathcal{W})$, and suppose $\mathcal{V}=\ker(T)\oplus\mathcal{M}$. Then $T|_{\mathcal{M}}:\mathcal{M}\to\operatorname{range}(T)$ is an isomorphism.
\begin{proof}
Indeed, $\ker(T|_{\mathcal{M}})=\mathcal{M}\cap\ker(T)=\{\mathbf{0}\}$, and every $T\mathbf{v}$ equals $T\mathbf{m}$ for the decomposition $\mathbf{v}=\mathbf{k}+\mathbf{m}$.\\
Thus rank-nullity counts $\mathcal{V}=\ker(T)\oplus\mathcal{M}$ with $\mathcal{M}\cong\operatorname{range}(T)$.
\end{proof}
\end{proposition}

\begin{definition} \hlt{Quotient Map}\\
Let $\mathcal{U}$ be a subspace of $\mathcal{V}$. The \hlt{quotient map} is the map $\pi : \mathcal{V} \to \mathcal{V}/\mathcal{U}$ defined by $\pi(\mathbf{v})=\mathbf{v}+\mathcal{U}$.
\end{definition}

\begin{proposition} \label{prop:quotient-map-properties} \hlt{Properties of the Quotient Map}\\
The quotient map $\pi$ is linear, surjective, and $\ker(\pi) = \mathcal{U}$.
\begin{proof}
For $\mathbf{v},\mathbf{w}\in\mathcal{V}$ and $a\in\mathbb{F}$, $\pi(a\mathbf{v}+\mathbf{w})=(a\mathbf{v}+\mathbf{w})+\mathcal{U} =a(\mathbf{v}+\mathcal{U})+(\mathbf{w}+\mathcal{U}) =a\pi(\mathbf{v})+\pi(\mathbf{w})$ so $\pi$ is linear. It is surjective because every coset in $\mathcal{V}/\mathcal{U}$ has the form $\mathbf{v}+\mathcal{U}=\pi(\mathbf{v})$. Finally, $\ker(\pi)=\{\mathbf{v}\in\mathcal{V}:\mathbf{v}+\mathcal{U}=\mathcal{U}\}=\mathcal{U}$.\\
\hlt{Structural Commentary}: The quotient map collapses $\mathcal{U}$ to zero and records cosets.\\
\hlt{Key Insight}: Quotient maps realise subspaces as kernels.
\end{proof}
\end{proposition}

\begin{theorem} \label{thm:quotient-factorisation-through-projection} \hlt{Universal Property of the Quotient Map}\\
Let $\mathcal{U}\leq\mathcal{V}$, $\pi:\mathcal{V}\to\mathcal{V}/\mathcal{U}$, and $T\in\LL(\mathcal{V},\mathcal{W})$ with $\mathcal{U}\leq\ker(T)$.\\
Then $\exists!\,\overline{T}\in\LL(\mathcal{V}/\mathcal{U},\mathcal{W})
$ such that $\overline{T}\pi=T$, $\overline{T}(\mathbf{v}+\mathcal{U})=T\mathbf{v}$, and $\ker(\overline{T})=\ker(T)/\mathcal{U}$, $\operatorname{range}(\overline{T})=\operatorname{range}(T)$.
\begin{proof}
Define $\overline{T}(\mathbf{v}+\mathcal{U})=T\mathbf{v}$. If $\mathbf{v}+\mathcal{U}=\mathbf{w}+\mathcal{U}$, then $\mathbf{v}-\mathbf{w}\in\mathcal{U}\leq\ker(T)$, so $T\mathbf{v}=T\mathbf{w}$; hence $\overline{T}$ is well-defined.\\
Also $\overline{T}\big(a(\mathbf{v}+\mathcal{U})+b(\mathbf{w}+\mathcal{U})\big) =T(a\mathbf{v}+b\mathbf{w}) =aT\mathbf{v}+bT\mathbf{w}$.\\
By definition, $\overline{T}\pi=T$. If $S\pi=T$, then $S(\mathbf{v}+\mathcal{U})=S\pi(\mathbf{v})=T\mathbf{v}=\overline{T}(\mathbf{v}+\mathcal{U})$, so $S=\overline{T}$.\\
Also $\mathbf{v}+\mathcal{U}\in\ker(\overline{T}) \Longleftrightarrow T\mathbf{v}=\mathbf{0} \Longleftrightarrow \mathbf{v}\in\ker(T)$.\\
So $\ker(\overline{T})=\ker(T)/\mathcal{U}$, and $\operatorname{range}(\overline{T})=\{\overline{T}(\mathbf{v}+\mathcal{U}):\mathbf{v}\in\mathcal{V}\}=\operatorname{range}(T)$.\\
\hlt{Structural Commentary}: Maps killing $\mathcal{U}$ factor through $\mathcal{V}/\mathcal{U}$.\\
\hlt{Key Insight}: Quotients are universal domains for maps with prescribed kernel ambiguity.
\end{proof}
\end{theorem}

\begin{theorem} \label{thm:subspace-correspondence-quotient} \hlt{Subspace Correspondence for Quotients}\\
Let $\pi:\mathcal{V}\to\mathcal{V}/\mathcal{U}$ be the quotient map. The assignments $\mathcal{T}\longmapsto \mathcal{T}/\mathcal{U}$ and $\mathcal{Y}\longmapsto \pi^{-1}(\mathcal{Y})$ are inverse inclusion-preserving bijections between $\{\mathcal{T}$ subspace of $\mathcal{V}:\mathcal{U}\subseteq\mathcal{T}\}$ and $\{\mathcal{Y}$ subspace of $\mathcal{V}/\mathcal{U}\}$.
\begin{proof}
If $\mathcal{U}\subseteq\mathcal{T}\subseteq\mathcal{V}$, then $\mathcal{T}/\mathcal{U}$ is a subspace of $\mathcal{V}/\mathcal{U}$.\\
If $\mathcal{Y}\subseteq\mathcal{V}/\mathcal{U}$ is a subspace, then $\pi^{-1}(\mathcal{Y})$ is a subspace of $\mathcal{V}$ and contains $\ker(\pi)=\mathcal{U}$.\\
For $\mathcal{U}\subseteq\mathcal{T}\subseteq\mathcal{V}$, $\pi^{-1}(\mathcal{T}/\mathcal{U})=\mathcal{T}$, since $\mathbf{v}+\mathcal{U}=\mathbf{t}+\mathcal{U}$ with $\mathbf{t}\in\mathcal{T}$ implies $\mathbf{v}-\mathbf{t}\in\mathcal{U}\subseteq\mathcal{T}$.\\
For $\mathcal{Y}\subseteq\mathcal{V}/\mathcal{U}$, surjectivity gives $\pi\left(\pi^{-1}(\mathcal{Y})\right)=\mathcal{Y}$, that is, $\pi^{-1}(\mathcal{Y})/\mathcal{U}=\mathcal{Y}$\\
Both assignments preserve inclusion by monotonicity of images and preimages.\\
\hlt{Structural Commentary}: Subspaces above $\mathcal{U}$ are exactly the subspaces visible after collapsing $\mathcal{U}$.\\
\hlt{Key Insight}: The quotient map translates containment above $\mathcal{U}$ into containment inside $\mathcal{V}/\mathcal{U}$.
\end{proof}
\end{theorem}

\begin{figure}[H]
  \centering
  \begin{tikzpicture}[
  x=0.88cm,
  y=0.88cm,
  >=Stealth,
  every node/.style={font=\small},
  point/.style={circle, fill=gray!75, inner sep=1.55pt},
  edge/.style={->, gray!65!black, line width=0.7pt},
  map/.style={->, gray!60, line width=0.65pt},
  hint/.style={font=\scriptsize, text=gray!70!black, inner sep=1pt}
]
  \node[point] (U) at (-2.4,0.55) {};
  \node[point] (T) at (0,0.55) {};
  \node[point] (V) at (2.4,0.55) {};

  \node[point] (QU) at (-2.4,-0.55) {};
  \node[point] (QT) at (0,-0.55) {};
  \node[point] (QV) at (2.4,-0.55) {};

  \draw[edge] (U) -- (T);
  \draw[edge] (T) -- (V);
  \draw[edge] (QU) -- (QT);
  \draw[edge] (QT) -- (QV);

  \draw[map] (U) -- (QU);
  \draw[map] (T) -- node[right, hint] {$\pi$} (QT);
  \draw[map] (V) -- (QV);

  \node at (-2.4,0.9) {$\mathcal{U}$};
  \node at (0,0.9) {$\mathcal{T}$};
  \node at (2.4,0.9) {$\mathcal{V}$};

  \node at (-2.4,-0.92) {$\{\mathcal{U}\}$};
  \node at (0,-0.92) {$\mathcal{T}/\mathcal{U}$};
  \node at (2.4,-0.92) {$\mathcal{V}/\mathcal{U}$};
\end{tikzpicture}

  \caption{Intermediate subspaces above $\mathcal{U}$ correspond to subspaces of the quotient.}
  \label{fig:quotient-correspondence-subspaces}
\end{figure}

\begin{theorem} \label{thm:quotient-by-kernel-is-range} \hlt{Quotient by the Kernel}\\
Let $T\in\LL(\mathcal{V},\mathcal{W})$. Then $\widetilde{T}(\mathbf{v}+\ker(T))\defeq T\mathbf{v}$ determines a well-defined linear map $\widetilde{T}:\mathcal{V}/\ker(T)\to\mathcal{W}$. Also,
\begin{enumerate}[label=(\roman*)]
\setlength{\itemsep}{0pt}
\item $\widetilde{T}$ is injective;
\item $\operatorname{range}(\widetilde{T})=\operatorname{range}(T)$;
\item $\mathcal{V}/\ker(T)\cong\operatorname{range}(T)$.
\end{enumerate}
\begin{proof}
If $\mathbf{v}+\ker(T)=\mathbf{v}'+\ker(T)$, then $\mathbf{v}-\mathbf{v}'\in\ker(T)$, so $T\mathbf{v}=T\mathbf{v}'$. Thus $\widetilde{T}$ is well-defined.\\
For cosets and $a\in\mathbb{F}$, $\widetilde{T}\big(a(\mathbf{v}+\ker(T))+(\mathbf{v}'+\ker(T))\big) =T(a\mathbf{v}+\mathbf{v}') =aT\mathbf{v}+T\mathbf{v}'$, so $\widetilde{T}$ is linear.
\begin{enumerate}[label=(\roman*)]
\setlength{\itemsep}{0pt}
\item If $\widetilde{T}(\mathbf{v}+\ker(T))=\mathbf{0}$, then $T\mathbf{v}=\mathbf{0}$, so $\mathbf{v}\in\ker(T)$ and $\mathbf{v}+\ker(T)=\ker(T)$, the zero coset.\\
Hence $\ker(\widetilde{T})=\{\ker(T)\}$ and $\widetilde{T}$ is injective.
\item For every coset $\mathbf{v}+\ker(T)$, the value $\widetilde{T}(\mathbf{v}+\ker(T))=T\mathbf{v}$ lies in $\operatorname{range}(T)$.\\
Conversely, every vector in $\operatorname{range}(T)$ has the form $T\mathbf{v}=\widetilde{T}(\mathbf{v}+\ker(T))$. Hence the ranges are equal.
\item Let $\widehat{T}:\mathcal{V}/\ker(T)\to\operatorname{range}(T)$ be the same formula as $\widetilde{T}$ with codomain restricted to $\operatorname{range}(T)$.\\
The linearity proved above and parts (i)--(ii) show that $\widehat{T}$ is linear, injective, and surjective, so Proposition \ref{prop:linear-map-invertible-iff-bijective} gives $\mathcal{V}/\ker(T)\cong\operatorname{range}(T)$.
\end{enumerate}
\hlt{Structural Commentary}: Quotienting by the kernel removes exactly the ambiguity of $T$.\\
\hlt{Key Insight}: Every linear map factors as a quotient followed by an isomorphism onto its range.
\end{proof}
\end{theorem}

\begin{proposition} \label{prop:quotient-by-complement-is-complement} \hlt{Quotient by a Complemented Subspace}\\
Let $\mathcal{U}$ and $\mathcal{W}$ be subspaces of $\mathcal{V}$ with $\mathcal{V}=\mathcal{U}\oplus\mathcal{W}$. Then $J:\mathcal{W}\to\mathcal{V}/\mathcal{U}$, $J(\mathbf{w})=\mathbf{w}+\mathcal{U}$ is an isomorphism.\\
Hence every complement of $\mathcal{U}$ is isomorphic to $\mathcal{V}/\mathcal{U}$.
\begin{proof}
The map is linear by the quotient operations from Proposition \ref{prop:quotient-space-structure}.\\
If $J(\mathbf{w})=\mathcal{U}$, then $\mathbf{w}\in\mathcal{U}\cap\mathcal{W}=\{\mathbf{0}\}$, so $J$ is injective.\\
If $\mathbf{v}+\mathcal{U}\in\mathcal{V}/\mathcal{U}$, write $\mathbf{v}=\mathbf{u}+\mathbf{w}$ with $\mathbf{u}\in\mathcal{U}$ and $\mathbf{w}\in\mathcal{W}$. Then $\mathbf{v}+\mathcal{U}=\mathbf{w}+\mathcal{U}=J(\mathbf{w})$, so $J$ is surjective.\\
\hlt{Structural Commentary}: A complement gives canonical representatives for quotient cosets.\\
\hlt{Key Insight}: Collapsing $\mathcal{U}$ leaves exactly the complementary directions visible.
\end{proof}
\end{proposition}

\begin{proposition} \label{prop:quotient-of-direct-sum-by-direct-sum} \hlt{Quotient of a Direct Sum}\\
If $\mathcal{V}=\mathcal{V}_1\oplus\mathcal{V}_2$ and $\mathcal{U}_j\leq\mathcal{V}_j$ for $j=1,2$, then $\frac{\mathcal{V}}{\mathcal{U}_1\oplus\mathcal{U}_2} \cong \frac{\mathcal{V}_1}{\mathcal{U}_1}\oplus \frac{\mathcal{V}_2}{\mathcal{U}_2}$.
\begin{proof}
Define, using the unique decomposition\\
$\mathbf{v}=\mathbf{v}_1+\mathbf{v}_2$,
$\Phi:\mathcal{V}\to (\mathcal{V}_1/\mathcal{U}_1)\oplus(\mathcal{V}_2/\mathcal{U}_2)$ and $\Phi(\mathbf{v}_1+\mathbf{v}_2) =(\mathbf{v}_1+\mathcal{U}_1,\mathbf{v}_2+\mathcal{U}_2)$.\\
For $\mathbf{x}=\mathbf{x}_1+\mathbf{x}_2$, $\mathbf{y}=\mathbf{y}_1+\mathbf{y}_2$, and $a,b\in\mathbb{F}$, $\Phi(a\mathbf{x}+b\mathbf{y}) =a\Phi(\mathbf{x})+b\Phi(\mathbf{y})$, so $\Phi$ is linear.\\
It is surjective since $(\mathbf{v}_1+\mathcal{U}_1,\mathbf{v}_2+\mathcal{U}_2)=\Phi(\mathbf{v}_1+\mathbf{v}_2)$. Finally, $\ker(\Phi)=\{\mathbf{v}_1+\mathbf{v}_2:\mathbf{v}_1\in\mathcal{U}_1,\ \mathbf{v}_2\in\mathcal{U}_2\} =\mathcal{U}_1\oplus\mathcal{U}_2$.\\
The result follows from Theorem \ref{thm:quotient-by-kernel-is-range}.\\
\hlt{Structural Commentary}: Quotienting respects direct-sum coordinates.\\
\hlt{Key Insight}: Direct-sum coordinates make quotienting componentwise.
\end{proof}
\end{proposition}

\begin{perspective} \label{per:quotients-and-complements-are-not-functorial-warnings} \hlt{Quotient and Complement Warnings}\\
If $\mathcal{V}=\mathcal{U}\oplus\mathcal{W}$, then $\mathcal{W}\cong\mathcal{V}/\mathcal{U}$, but complements are not functorial in $\mathcal{U}$.\\
For $\mathcal{V}$ countably infinite-dimensional with basis $(\mathbf{e}_n)_{n\geq1}$, put $\mathcal{S}=\operatorname{span}(\mathbf{e}_n:n\geq2)$.\\
Then $\mathcal{S}\cong\mathcal{V}$ and $\mathcal{V}=\mathcal{V}\oplus\{\mathbf{0}\} =\mathcal{S}\oplus\mathbb{F}\mathbf{e}_1$; $\mathcal{V}/\mathcal{V}=\{\mathbf{0}\}$, $\mathcal{V}/\mathcal{S}\cong\mathbb{F}\mathbf{e}_1$.
\end{perspective}

\begin{definition} \label{def:codimension-subspace} \hlt{Codimension}\\
If $\mathcal{U}$ is a subspace of $\mathcal{V}$ and $\dim(\mathcal{V}/\mathcal{U})<\infty$, define $\operatorname{codim}_{\mathcal{V}}(\mathcal{U})\defeq\dim(\mathcal{V}/\mathcal{U})$.\\
Thus, if $\mathcal{V}$ is finite-dimensional, Theorem \ref{thm:dimension-of-quotient} gives $\operatorname{codim}_{\mathcal{V}}(\mathcal{U})=\dim(\mathcal{V})-\dim(\mathcal{U})$.\\
If $\mathcal{V}=\mathcal{U}\oplus\mathcal{W}$, Proposition \ref{prop:quotient-by-complement-is-complement} gives $\operatorname{codim}_{\mathcal{V}}(\mathcal{U})=\dim(\mathcal{W})$.
\end{definition}

\begin{perspective} \label{per:quotients-split-over-fields} \hlt{Quotients Split Over Fields}\\
In finite-dimensional vector spaces, choosing a complement $\mathcal{V}=\mathcal{U}\oplus\mathcal{W}$ identifies $\mathcal{V}/\mathcal{U}$ with $\mathcal{W}$. The quotient map collapses $\mathcal{U}$ and keeps the complementary directions.
\end{perspective}

\begin{proposition} \label{prop:finite-dimensional-one-sided-inverses-split} \hlt{One-Sided Inverses Split in Finite Dimensions}\\
Let $\mathcal{V}$ and $\mathcal{W}$ be finite-dimensional vector spaces and let $T\in\LL(\mathcal{V},\mathcal{W})$.
\begin{enumerate}[label=(\roman*)]
\setlength{\itemsep}{0pt}
\item If $T$ is injective, then there exists $L\in\LL(\mathcal{W},\mathcal{V})$ such that $LT=I_{\mathcal{V}}$.
\item If $T$ is surjective, then there exists $R\in\LL(\mathcal{W},\mathcal{V})$ such that $TR=I_{\mathcal{W}}$.
\end{enumerate}
\begin{proof}
\begin{enumerate}[label=(\roman*)]
\setlength{\itemsep}{0pt}
\item Choose $\mathcal{H}\subseteq\mathcal{W}$ with $\mathcal{W}=\operatorname{range}(T)\oplus\mathcal{H}$. For $\mathbf{w}=T\mathbf{v}+\mathbf{h}$, define $L\mathbf{w}=\mathbf{v}$. This is well-defined since the direct-sum decomposition is unique and $T$ is injective. It is linear, and $LT\mathbf{v}=\mathbf{v}$ for every $\mathbf{v}\in\mathcal{V}$.
\item Choose $\mathcal{H}\subseteq\mathcal{V}$ with $\mathcal{V}=\ker(T)\oplus\mathcal{H}$. Then $T|_{\mathcal{H}}:\mathcal{H}\to\mathcal{W}$ is injective since $\mathcal{H}\cap\ker(T)=\{\mathbf{0}\}$, and surjective since each $\mathbf{v}=\mathbf{k}+\mathbf{h}$ satisfies $T\mathbf{v}=T\mathbf{h}$. Let $R=(T|_{\mathcal{H}})^{-1}$, viewed as a map $\mathcal{W}\to\mathcal{V}$.\\
Then $R$ is linear and $TR=I_{\mathcal{W}}$.
\end{enumerate}
\hlt{Structural Commentary}: Complements supply the missing directions on which a one-sided inverse is chosen.\\
\hlt{Key Insight}: Finite-dimensional vector spaces split kernels and ranges.
\end{proof}
\end{proposition}

\begin{proposition} \label{prop:isomorphisms-preserve-bases} \hlt{Isomorphisms Preserve Bases}\\
Let $T:\mathcal{V}\to\mathcal{W}$ be an isomorphism and let $\beta=(\mathbf{v}_1,\ldots,\mathbf{v}_n)$ be a finite list in $\mathcal{V}$. Set $T\beta=(T\mathbf{v}_1,\ldots,T\mathbf{v}_n)$.\\
The following are equivalent:
\begin{enumerate}[label=(\roman*)]
\setlength{\itemsep}{0pt}
\item $\beta$ is a basis of $\mathcal{V}$.
\item $T\beta$ is a basis of $\mathcal{W}$.
\end{enumerate}
\begin{proof}
\begin{enumerate}[label=(\roman*).]
\setlength{\itemsep}{0pt}
\item[$(i)$]$\Rightarrow (ii)$ If $\mathbf{w}\in\mathcal{W}$, choose $\mathbf{v}\in\mathcal{V}$ with $\mathbf{w}=T\mathbf{v}$ and write $\mathbf{v}=\sum_j a_j\mathbf{v}_j$. Then $\mathbf{w}=\sum_j a_jT\mathbf{v}_j$, so $T\beta$ spans $\mathcal{W}$. If $\sum_j a_jT\mathbf{v}_j=\mathbf{0}$, applying $T^{-1}$ gives $\sum_j a_j\mathbf{v}_j=\mathbf{0}$; hence all $a_j=0$. Thus $T\beta$ is a basis.
\item[$(ii)$]$\Rightarrow (i)$
Assume $T\beta$ is a basis of $\mathcal{W}$. If $\mathbf{v}\in\mathcal{V}$, write $T\mathbf{v}=\sum_j a_jT\mathbf{v}_j$; applying $T^{-1}$ gives $\mathbf{v}=\sum_j a_j\mathbf{v}_j$, so $\beta$ spans $\mathcal{V}$. If $\sum_j a_j\mathbf{v}_j=\mathbf{0}$, applying $T$ gives $\sum_j a_jT\mathbf{v}_j=\mathbf{0}$; hence all $a_j=0$. Thus $\beta$ is linearly independent, and therefore a basis of $\mathcal{V}$.
\end{enumerate}
\hlt{Structural Commentary}: Isomorphisms preserve exactly the spanning and independence data that define a basis.\\
\hlt{Key Insight}: An isomorphism changes the ambient space, not the linear relations among vectors.
\end{proof}
\end{proposition}

\begin{theorem} \label{thm:finite-dimensional-operator-invertible-equivalences} \hlt{Finite-Dimensional Invertibility Criteria}\\
If $\mathcal{V}$ is finite-dimensional and $T \in \LL(\mathcal{V})$, then the following are equivalent:
\begin{enumerate}[label=\roman*.]
\setlength{\itemsep}{0pt}
\item $T$ is invertible;
\item $T$ is injective;
\item $T$ is surjective.
\end{enumerate}
\begin{proof}
\begin{enumerate}[label=(\roman)*.]
\setlength{\itemsep}{0pt}
\item[$(i)$] $\Rightarrow (ii)$ follows from Proposition \ref{prop:linear-map-invertible-iff-bijective}.
\item[$(ii)$] $\Rightarrow (iii)$: If $T$ is injective, then $\ker(T)=\{\mathbf{0}\}$.\\
Rank-nullity gives $\dim(\operatorname{range}(T))=\dim(\mathcal{V})$, hence $\operatorname{range}(T)=\mathcal{V}$, so $T$ is surjective.
\item[$(iii)$] $\Rightarrow (i)$: If $T$ is surjective, rank-nullity gives $\dim(\ker(T))=0$, so $T$ is injective.\\
Thus $T$ is bijective, hence invertible by Proposition \ref{prop:linear-map-invertible-iff-bijective}.
\end{enumerate}
\hlt{Structural Commentary}: For operators on finite-dimensional spaces, kernel and range dimensions trade off exactly. Injectivity and surjectivity therefore force each other.\\
\hlt{Key Insight}: In equal finite dimensions, losing no information is the same as reaching the whole space.
\end{proof}
\end{theorem}

\begin{perspective} \label{per:kernel-obstruction-to-invertibility} \hlt{Kernel Obstruction to Invertibility}\\
By Theorem \ref{thm:finite-dimensional-operator-invertible-equivalences} and Proposition \ref{prop:injective-iff-kernel-zero}, a finite-dimensional operator is invertible exactly when its kernel is zero. Any nonzero $\mathbf{v}\in\ker T$ therefore blocks invertibility.
\end{perspective}

\begin{theorem} \label{thm:finite-dimensional-isomorphic-iff-same-dimension} \hlt{Isomorphism and Dimension}\\
Two finite-dimensional vector spaces are isomorphic if and only if they have the same dimension.
\begin{proof}
\begin{enumerate}[label=\roman*.]
\setlength{\itemsep}{0pt}
\item[$(\Rightarrow)$] If $T:\mathcal{V}\to\mathcal{W}$ is an isomorphism and $(\mathbf{v}_1,\ldots,\mathbf{v}_n)$ is a basis of $\mathcal{V}$, then Proposition \ref{prop:isomorphisms-preserve-bases} shows that $(T\mathbf{v}_1,\ldots,T\mathbf{v}_n)$ is a basis of $\mathcal{W}$.
\item[$(\Leftarrow)$] If both spaces have dimension $n$, choose bases $(\mathbf{v}_1,\ldots,\mathbf{v}_n)$ and $(\mathbf{w}_1,\ldots,\mathbf{w}_n)$.\\
By Proposition \ref{prop:linear-map-determined-by-basis}, there is a linear map $T$ with $T\mathbf{v}_j=\mathbf{w}_j$ for every $j$. Then $T(\sum a_j\mathbf{v}_j)=\sum a_j\mathbf{w}_j$, so $\ker(T)=\{\mathbf{0}\}$ and $\operatorname{range}(T)=\mathcal{W}$. Hence $T$ is an isomorphism.
\end{enumerate}
\hlt{Structural Commentary}: Dimension classifies finite-dimensional spaces up to isomorphism.\\
\hlt{Key Insight}: Matching bases give an isomorphism; any isomorphism carries bases to bases.
\end{proof}
\end{theorem}

